% !TeX document-id = {ceb9b0f7-fe16-4e6e-808c-436a33f8cebf}
% !TeX program = pdflatex
% !TeX TXS-program:compile = txs:///pdflatex/[-shell-escape]
% !TeX TXS-program:pdflatex = pdflatex -synctex=1 -interaction=nonstopmode --shell-escape %.tex

\documentclass[11pt,a4paper,oneside]{report}

\usepackage{style/cybersec}
\usepackage{subfiles}

\begin{document}

% ================================================================================
% Portada Principal del Libro Recopilatorio (Geometric Blueprint - Fondo Claro)
% ================================================================================
\begin{titlepage}
	\begin{tikzpicture}[remember picture,overlay]
		% Esquina superior derecha: Acentos geométricos
		\fill[csecNavy] (current page.north east) -- ([xshift=-6.5cm]current page.north east) -- ([yshift=-4.2cm]current page.north east) -- cycle;
		\fill[csecCyan] ([xshift=-6.5cm]current page.north east) -- ([xshift=-6.9cm]current page.north east) -- ([yshift=-4.5cm]current page.north east) -- ([yshift=-4.2cm]current page.north east) -- cycle;
		% Esquina inferior izquierda: Acentos geométricos
		\fill[csecNavy] (current page.south west) -- ([xshift=6.5cm]current page.south west) -- ([yshift=4.2cm]current page.south west) -- cycle;
		\fill[csecCyan] ([xshift=6.5cm]current page.south west) -- ([xshift=6.9cm]current page.south west) -- ([yshift=4.5cm]current page.south west) -- ([yshift=4.2cm]current page.south west) -- cycle;
		% Marco perimetral fino
		\draw[csecBorder, line width=0.6pt] ([xshift=1.2cm, yshift=-1.2cm]current page.north west) rectangle ([xshift=-1.2cm, yshift=1.2cm]current page.south east);
	\end{tikzpicture}

	\centering
	\vspace*{2.8cm}

	% Título Principal
	{\fontsize{34}{40}\selectfont \textbf{\color{csecDark}CYBERSECURITY}}\par
	\vspace{0.2cm}
	{\fontsize{26}{32}\selectfont \color{csecCyan}\textbf{PLAYBOOK \& WRITEUPS}}\par
	\vspace{1.8cm}

	% Colección de Plataformas y Entornos Evaluados (Fácilmente extensible)
	\begin{tcolorbox}[
			enhanced,
			colback=csecLight,
			colframe=csecBorder,
			boxrule=0.8pt,
			arc=5pt,
			width=0.85\textwidth,
			top=14pt, bottom=14pt, left=14pt, right=14pt
		]
		\centering
		{\normalsize \textbf{\color{csecDark}Plataformas y Entornos Evaluados}}\par
		\vspace{12pt}
		\platformBadge{HackTheBox} \hspace{4pt}
		\platformBadge{TryHackMe} \hspace{4pt}
		\platformBadge{PortSwigger} \hspace{4pt}
		\platformBadge{pwn.college} \hspace{4pt}
		\platformBadge{CTF}
	\end{tcolorbox}

	\vfill

	% Pie de Portada: Autor, Versión y Fecha
	\begin{minipage}{0.75\textwidth}
		\centering
		{\Large \textbf{\color{csecDark}errantprogrammer}}\par
		\vspace{0.3cm}
		{\footnotesize \color{csecCyan}\textbf{v1.0 (2026)} \quad \color{csecBorder}\textbar\quad \color{csecMuted}\today}
	\end{minipage}
	\vspace*{1.2cm}
\end{titlepage}

% ================================================================================
% Páginas Preliminares (Dashboard e Índice de Contenidos)
% ================================================================================
\pagenumbering{roman}
\pagestyle{fancy}
\setCustomHeader{{\color{csecCyan}\faChartBar}\ \textbf{Playbook}}{Dashboard de Progreso}{\today}

\section*{Dashboard de Máquinas y Desafíos Resueltos}
\addcontentsline{toc}{chapter}{Dashboard de Progreso}

\renderDashboard

\vspace{0.5cm}

% ================================================================================
% Tabla de Contenidos General
% ================================================================================
\cleardoublepage
\setCustomHeader{{\color{csecCyan}\faList}\ \textbf{Playbook}}{Índice General}{\today}
\tableofcontents
\cleardoublepage

% ================================================================================
% Inicio de Capítulos Técnicos
% ================================================================================
\pagenumbering{arabic}

% ================================================================================
% CAPÍTULO 1: GUÍA DE USO Y MANUAL TÉCNICO
% ================================================================================
\subfile{chapters/guide/template_guide}

% ================================================================================
% CAPÍTULO 2: KNOWLEDGE BASE & BITÁCORA DE APRENDIZAJE
% ================================================================================
\subfile{chapters/knowledge_base/knowledge_base}

% % ================================================================================
% % PARTE I: HACK THE BOX
% % ================================================================================
% \chapter{Hack The Box}
% \label{chap:hackthebox}
% \setPlatformHeader{HackTheBox}{Hack The Box}{\today}
% En este capítulo se recopilan los informes técnicos detallados de las máquinas resueltas en la plataforma \textbf{Hack The Box}.
%
% % Inclusión modular del writeup de Lame
%
% % ================================================================================
% % PARTE II: TRYHACKME
% % ================================================================================
% \chapter{TryHackMe}
% \label{chap:tryhackme}
% \setPlatformHeader{TryHackMe}{TryHackMe}{\today}
% En este capítulo se documentan las salas temáticas y redes de práctica completadas en la plataforma \textbf{TryHackMe}.
%
% \begin{notebox}[Salas en Progreso]
% 	Para incorporar una nueva sala de TryHackMe, utiliza la plantilla ubicada en \texttt{writeups/tryhackme/\_template\_thm.tex} y enlázala en este capítulo mediante el comando \texttt{\textbackslash subfile\{writeups/tryhackme/nombre\_sala/writeup\}}.
% \end{notebox}
%
% % ================================================================================
% % PARTE III: PORTSWIGGER WEB SECURITY ACADEMY
% % ================================================================================
% \chapter{PortSwigger Web Security Academy}
% \label{chap:portswigger}
% \setPlatformHeader{PortSwigger}{PortSwigger Academy}{\today}
% Documentación y soluciones de laboratorios centrados en vulnerabilidades de aplicaciones web (OWASP Top 10, SQLi, XSS, CSRF, SSRF, Deserialización, etc.).
%
% \begin{notebox}[Laboratorios Web]
% 	Utiliza la plantilla en \texttt{writeups/portswigger/\_template\_portswigger.tex} para documentar vulnerabilidades web individuales.
% \end{notebox}
%
% % ================================================================================
% % PARTE IV: CTF CHALLENGES & COMPETENCIAS
% % ================================================================================
% \chapter{CTF Challenges \& Competencias}
% \label{chap:ctfs}
% \setPlatformHeader{CTF}{CTF Challenges}{\today}
% Resolución de retos específicos clasificados por categorías: Criptografía, Ingeniería Inversa, Explotación de Binarios (Pwn), Análisis Forense y OSINT.


\end{document}
